\Subhead{3.8.2}{Дифракция на случайном фазовом экране. Тепловое электромагнитное поле. Теорема взаимности}

Тонкий случайный фазовый экран преобразует регулярное поле как $E_+=E_-\exp[\ii\phi(\bm\rho)]$. Если фаза гауссовская с постоянным средним, то
\begin{equation}
 \langle\e^{\ii[\phi(\bm\rho_1)-\phi(\bm\rho_2)]}\rangle
 =\exp[-D_\phi(\bm\rho_1,\bm\rho_2)/2],
\end{equation}
где $D_\phi=\langle(\phi_1-\phi_2)^2\rangle$. Структурная функция фазы непосредственно уменьшает взаимную когерентность. Последующее распространение переводит фазовые флуктуации в амплитудные: возникают мерцания и спекл. Среднее поле и средняя интенсивность различны; при случайной фазе первое может исчезнуть при ненулевой второй.

Равновесное тепловое поле близко к нулевому круговому гауссовскому полю, если складывается много независимых мод. Для одной поляризации и выбранного пространственного канала спаривание даёт соотношение Зигерта $g^{(2)}(\tau)=1+|g^{(1)}(\tau)|^2$. При $\tau=0$ идеальное тепловое поле имеет $g^{(2)}(0)=2$; усреднение по независимым модам уменьшает контраст.

Теорема Лоренца о взаимности для линейной среды с симметричными тензорами $\varepsilon,\mu$ следует из вычитания двух тождеств Пойнтинга:
\begin{equation}
 \nabla\cdot(\bm E_1\times\bm H_2-\bm E_2\times\bm H_1)
 =\bm E_2\cdot\bm J_1-\bm E_1\cdot\bm J_2.
\end{equation}
При исчезающем поверхностном вкладе интегралы взаимного возбуждения равны; функция Грина симметрична с перестановкой точек и индексов. Потери сами по себе не нарушают взаимность, если тензоры остаются симметричными. Магнитоактивная среда может быть невзаимной, поэтому переставлять источник и приёмник без условий нельзя.
